\chapter{Flatness}
\label{chap:flatness}
Recall from \Cref{def: flat} that we defined an $A$-module $M$ to be flat if tensoring by $M$ preserves exactness. As tensor is already right-exact, this is equivalent to preserving injection, i.e., if $L\to N$ is injective, then so is the induced map $L\otimes_A M\to N\otimes_A M$. In this chapter, we will explore the properties and uses of flat modules.

The first thing that you probably have noticed is that
in order to prove that $M$ is flat, one needs to check
that it preserves injectivity for all
injective morphisms $L\to N$ (including infintely-generated ones). 
This is a horrible condition to check.
This, in \Cref{sec:homological_algebra} and \Cref{sec:flatness_and_tor}, we develop
conditions to check if a module is flat (or not flat).

The cleanest way to do this require formidable machinery
in homological algebra, which borrows tools from algebraic
topology to study algebra. One can motivate this as follows:
whenever you see a right exact functor, like
$0\to L\to N\to P\to 0$, which implies the right exact sequence
$$L\otimes_A M \longto N\otimes_A M \longto P\otimes_A M\longto 0,$$
you should dream of extending the sequence indefinitely leftward,
giving the \vocab{left-derived functor}.
In this case, the exact sequence we will construct is
    \begin{center}
    \begin{tikzcd}
    & \dots \arrow[r]
    & \Tor_2^A(P,M) \arrow[lld] \\
    \Tor_1^A(L,M)  \arrow[r]
    & \Tor_1^A(N,M) \arrow[r]
    & \Tor_1^A(P,M) \arrow[lld] \\
    L\otimes_A M  \arrow[r]
    & N\otimes_A M \arrow[r]
    & P\otimes_A M \arrow[r] & 0
    \end{tikzcd}
\end{center}
In some sense, $\Tor$ measures the failure of exactness
of tensor product, and thus useful for
understanding flatness.
We will define $\Tor_i^A(L,M)$ at the beginning of 
\Cref{sec:flatness_and_tor}.
Then, we provide various flatness criteria
that we get from understanding the $\Tor$ functor.

In algebraic geometry, flat modules are tied closely
with flat morphisms.
\begin{definition}
    A ring morphisms $\phi : A\to B$ is \vocab[flat morphisms]{flat}
    if and only if it promotes $B$ to a flat $A$-algebra.
\end{definition}
Geometrically, flat morphism means ``nicely varying''.
An example of morphisms that are not nicely varying is
\begin{align*}
\phi : \CC[x] &\to \frac{\CC[x,y]}{(xy)} \\
x &\mapsto x
\end{align*}
Geometrically, this corresponds to the map from
variety $\{xy=0\}$ to an affine line $\{x\in\CC\}$ by projecting to
the $x$ coordinate. One can check that
$\CC[x,y]/(xy)$ is not a flat $\CC[x]$-algebra
(we leave this as \Cref{prob:not_flat_example}).
This is not ``nicely-varying'' because
the map is injective at all points, except above $x=0$,
where an entire line is squished into a single point.
The key result of this sort is flatness implies
constant dimension of fiber, which we prove in 
\Cref{sec:flat_const_dim}.

In contrast, the map 
\begin{align*}
\phi : \CC[x] &\to \frac{\CC[x,y]}{(xy-1)} \\
x &\mapsto x
\end{align*}
is flat (by \Cref{cor:localization_flat}
that localization is a flat module), 
even though we have seen that it does not look 
like a covering map. 
These two conditions are very different,
and flatness is more useful for maps that \emph{collapse
the dimension}.
\begin{center}
\begin{adjustbox}{valign=m}
\begin{asy}
size(4cm,0);
draw((-1,1)--(1,1), red+linewidth(1));
draw((0,1.5)--(0,0.5), red+linewidth(1));
label("\large $\Downarrow$", (0,0));
draw((-1,-0.5)--(1,-0.5), red+linewidth(1));
dot((0,1));
label("not flat", (0,-1),S);
\end{asy} 
\end{adjustbox}
\hspace{2cm}
\begin{adjustbox}{valign=m}
\begin{asy}
size(4cm,0);
import graph;
real f(real x){ return 0.1/x;}
draw(shift(0,1.5)*graph(f,0.1,1), red+linewidth(1));
draw(shift(0,1.5)*graph(f,-0.1,-1), red+linewidth(1));
label("\large $\Downarrow$", (0,0));
draw((0,2.5)--(0,0.5), dashed);
draw((-1,1.5)--(1,1.5), dashed);
draw((-1,-0.5)--(1,-0.5), red+linewidth(1));
dot((0,1.5));
label("flat", (0,-1),S);
\end{asy}
\end{adjustbox}
\end{center}

\section{Homological Algebra Crash Course}
\label{sec:homological_algebra}
This section provides the relevant background in homological 
algebra.
\subsection{Homology}
\paragraph{Chain Complexes}
\begin{definition}
    A \vocab{chain complex} is a sequence \[\cdots\stackrel{\partial_3}\longto A_2\stackrel{\partial_2}\longto A_1\stackrel{\partial_1}\longto A_0\stackrel{\partial_0}\longto A_{-1}\stackrel{\partial_{-1}}\longto A_{-2}\stackrel{\partial_{-2}}\longto\cdots\] such that $\partial_{i-1}\circ \partial_i = 0$ for all $i$ (i.e., $\Im \partial_i \subseteq \Ker \partial_{i-1}$).

    (Most chain complexes we consider will be finite. Strictly speaking, they can be extended infinitely in both directions with $0$'s.)
    \notation{$A_*$}{a chain complex}
\end{definition}
\paragraph{Homology}
\begin{definition}
    Given a chain complex $A_*$ as above, the \vocab{$n$-th homology group} is defined as 
    \notation{$H_n(A_*)$}{$n$-th homology group}
    \[H_n(A_*) := \Ker \partial_n / \Im \partial_{n+1}.\]
\end{definition}
\begin{example}
    Any exact sequence has all of its homology groups equal to $0$. Thus, the homology groups can be thought of as how ``far from exact'' a chain complex is.
\end{example}
\begin{theorem}[A variant of Snake lemma]
    \label{thm:long_exact_homology}
    Given a short exact sequence 
    $0\longto A_*\longto B_*\longto C_*\longto 0$
    of chain complex, i.e., the commuting diagram
\begin{center}
\begin{tikzcd}
    & 0 \arrow[d] & 0 \arrow[d] & 0 \arrow[d] & \\
    \dots \arrow[r]
    & A_{n+1} \arrow[r,"\partial"] \arrow[d]
    & A_n \arrow[r,"\partial"] \arrow[d]
    & A_{n-1} \arrow[r] \arrow[d]
    & \dots \\ 
    \dots \arrow[r]
    & B_{n+1} \arrow[r,"\partial"] \arrow[d]
    & B_n \arrow[r,"\partial"] \arrow[d]
    & B_{n-1} \arrow[r] \arrow[d]
    & \dots \\

    \dots \arrow[r]
    & C_{n+1} \arrow[r,"\partial"] \arrow[d]
    & C_n \arrow[r,"\partial"] \arrow[d]
    & C_{n-1} \arrow[r] \arrow[d]
    & \dots \\
    & 0 & 0 & 0 & 
\end{tikzcd}.
\end{center}
Then, there exists an exact sequence of homology groups
\begin{center}
    \begin{tikzcd}
	\dots \arrow[r]
	& H_{n+1}(B_*) \arrow[r] 
	& H_{n+1}(C_*) \arrow[lld, "\delta", swap] \\
	H_n(A_*) \arrow[r, swap] 
	& H_n(B_*) \arrow[r] 
	& H_n(C_*) \arrow[lld, "\delta", swap] \\
	H_{n-1}(A_*) \arrow[r, swap] 
	& H_{n-1}(B_*) \arrow[r]
	& \dots
    \end{tikzcd},
    \end{center}
\end{theorem}
This is a more general version of Snake Lemma
(\Cref{lem:snake}) because if
$A_*$, $B_*$, and $C_*$ each have two nonzero modules,
then we recover \Cref{lem:snake}.
It can be proven by a similary tedious argument,
so we omit the proof.
\subsection{Free Resolutions}
\begin{definition}[Free Resolution]
    Given an $A$-module $M$, a \vocab{free resolution} of $M$ is a chain complex of free $A$-modules which has $0$-th homology group $M$ and $n$-th homology group $0$ for $n\neq 0$.
\end{definition}
\begin{example}
    As a $\ZZ$-module, $\ZZ/n\ZZ$ has a free resolution given by \[\cdots \longto 0 \longto \ZZ \stackrel{\cdot n}\longto \ZZ \longto 0 \longto \cdots.\]
\end{example}
Every module $M$ has a free resolution. Indeed, we can first take a free module $M_0$ with a surjection $\varphi:M_0\to M$. Then, take a free module $M_1$ with a surjection $\partial_1: M_1\to \Ker \varphi$. Take a free module $M_2$ with a surjection $\partial_2:M_2\to \Ker \partial_1$, and so on. The resulting complex \[\cdots\stackrel{\partial_3}\longto M_2\stackrel{\partial_2}\longto M_1\stackrel{\partial_1}\longto M_0\longto 0\longto\cdots\] is a free resolution of $M$ by construction. (The process may not terminate, but that's fine because free resolutions need not be finite.)

\section{Flatness and Tor}
\label{sec:flatness_and_tor}
\subsection{The Tor Functor}
Every right-exact functor induces a left derived functor. In the case of the tensor product, the derived functor is $\Tor$.
\begin{definition}[The $\Tor$ functor]
    \label{def:tor}
    \notation{$\Tor_n^A(M,N)$}{$n$-th Tor functor}
    Given $A$-modules $M$ and $N$, let $N_*$ be a free resolution of $N$, and let $M\otimes N_*$ denote the chain complex induced by tensoring each component of $N_*$ with $M$. Then, the \vocab{$n$-th Tor functor} is defined as \[\Tor_n^A(M,N) := H_n(M\otimes N_*).\]
\end{definition}
In particular, $\Tor_0^A$ is simply $\otimes_A$.
\begin{example}[$\Tor$ of $\ZZ$-modules]
    We will compute $\Tor_i^\ZZ(\ZZ/m\ZZ, \ZZ/n\ZZ)$. Using the free resolution of $\ZZ/n\ZZ$ in the previous section, we simply want the homology groups of \[\cdots\to 0\longto \ZZ/m\ZZ\stackrel{\cdot n}\longto \ZZ/m\ZZ \longto 0\longto\cdots.\] The image of multiplication by $n$ in $\ZZ/m\ZZ$ is all the multiples of $\gcd(m,n)$, and the kernel is isomorphic to $\ZZ/\gcd(m,n)\ZZ$, so \[\Tor_0^\ZZ(\ZZ/m\ZZ,\ZZ/n\ZZ) \simeq \Tor_1^\ZZ(\ZZ/m\ZZ,\ZZ/n\ZZ) \simeq \ZZ/\gcd(m,n)\ZZ,\] and $\Tor_i^\ZZ(\ZZ/m\ZZ,\ZZ/n\ZZ)=0$ for $i\ge 2$.
\end{example}
Tor satisfies some nice properties, but proving many them is beyond the scope of this chapter.
\begin{theorem}[Properties of $\Tor$]
\label{thm:tor_properties}
\begin{alphenum}
    \item $\Tor_n^A(M,N)$ does not depend on the choice of free resolution $N_*$.
    \item If $0\to L\to N\to P\to 0$ is a short exact sequence of $A$-modules, then there is a long exact sequence 
    \begin{center}
    \begin{tikzcd}
    & \dots \arrow[r]
    & \Tor_2^A(P,M) \arrow[lld] \\
    \Tor_1^A(L,M)  \arrow[r]
    & \Tor_1^A(N,M) \arrow[r]
    & \Tor_1^A(P,M) \arrow[lld] \\
    L\otimes_A M  \arrow[r]
    & N\otimes_A M \arrow[r]
    & P\otimes_A M \arrow[r] & 0
    \end{tikzcd}
    \end{center}
    \item $\Tor_n^A(M,N) \simeq \Tor_n^A(N,M)$.
    \item If $I$ is a directed set and $(N_i)_{i\in I}$
    is a diagram, then
    $$\Tor_n^A\left(M, \varinjlim_{i\in I} N_i\right)
    = \varinjlim_{i\in I} \Tor_n^A(M, N_i)$$
    (cf. colimit in \Cref{sec:colimit}).
\end{alphenum}
\end{theorem}
\begin{proof}
\begin{alphenum}
    \item The proof of this is elementary,
    but it requires slightly more homological algebra 
    than what we have given. Thus, this is deferred to
    \Cref{thm:tor_free_res}.
    \item Take a free resolution 
    $\dots \to M_2\to M_1\to M_0$ of $M$. 
    Since $M_i$ is free (hence flat), we get 
    a commutative diagram with exact rows
    \begin{center}
    \begin{tikzcd}
        & \dots \arrow[d] & \dots \arrow[d] & \dots \arrow[d] \\
        0 \arrow[r] 
        & L\otimes_A M_i  \arrow[r] \arrow[d]
        & N\otimes_A M_i  \arrow[r] \arrow[d]
        & P\otimes_A M_i  \arrow[r]\arrow[d]
        & 0 \\
        0 \arrow[r] 
        & L\otimes_A M_{i-1}  \arrow[r] \arrow[d]
        & N\otimes_A M_{i-1}  \arrow[r] \arrow[d]
        & P\otimes_A M_{i-1}  \arrow[r] \arrow[d]
        & 0  \\
        & \dots & \dots & \dots
    \end{tikzcd}
    \end{center}
    In particular, we have the exact sequence
    of chain complexes 
    $0\to L\otimes_A M_* \to N\otimes_A M_*
    \to P\otimes_A M_*\to 0$.
    Then, the homology long exact sequence
    \Cref{thm:long_exact_homology} gives the desired
    long exact sequence.
    \item Intuitively, this should be true because
    tensor product is symmetric ($M\otimes_A N = N\otimes_A M$),
    but this turns out to be a very difficult one to prove.
    Roughly speaking, we take a free resolution
    of both $M$ and $N$. Then, take their product
    $M_i\otimes_A N_j$ to get a double complex 
    (that looks like an infinite grid). 
    Then, one manipulates
    homology groups of the double complex and proves the result
    (typically, this is done via spectral sequence).
    We refer the reader to \cite[Thm. 2.7.2]{weibel}.
    \item This follows from the following two facts:
    \begin{itemize}
    \item Colimit preserves exactness (\Cref{thm:direct_lim_is_exact}), 
    hence preserving kernels, images, and homology
    groups in a chain complex.
    \item Colimit commutes with tensor product:
    $M\otimes_A \varinjlim_{i\in I} N_i = 
    \varinjlim_{i\in I} (M\otimes_A N_i)$
    by \Cref{prob:colimit_tensor}.
    \end{itemize}
    Thus, every step in the definition of $\Tor$
    commutes with the colimit, so $\Tor$ commutes
    with colimit. \qedhere
\end{alphenum}
\end{proof}

% TODO mark help
\subsection{Ideal-theoretic Criterion of Flatness}
\begin{theorem}[Ideal-theoretic Criterion of Flatness]
    \label{thm:ideal_flat}
    Let $M$ be an $A$-module. The following are equivalent:
    \begin{alphenum}
        \item $M$ is flat.
        \item $\Tor_i^A(M,N) = 0$ for all $A$-modules $N$ and $i\ge 1$.
        \item $\Tor_1^A(M,N)=0$ for all $A$-modules $N$.
        \item $\Tor_1^A(M, A/I) = 0$ for all ideals $I\subseteq A$.
        %\item $\Tor_1^A(M, A/I) = 0$ for all finitely-generated ideals $I\subseteq A$.
    \end{alphenum}
\end{theorem}
\begin{proof}
    Note that $\text{(b)}\Rightarrow\text{(c)}\Rightarrow \text{(d)}$ is trivial.

    \textbf{(a) $\Rightarrow$ (b).} Consider a free resolution $N_*$ of $N$ given by \[\cdots \longto N_2\longto N_1\longto N_0\longto \cdots,\] and consider the long exact sequence \[\cdots \longto N_2\longto N_1 \longto N_0.\] As $M$ is flat, tensoring by $M$ preserves exactness, so \[\cdots \longto M\otimes N_2\longto M\otimes N_1\longto M\otimes N_0\] is also exact. Thus, the $i$-th homology group of $M\otimes N_*$ is $0$ for all $i\ge 1$.

    %\textbf{(e) $\Rightarrow$ (d).}
    
    \textbf{(d) $\Rightarrow$ (a).} 
    We prove this in three steps.
    \begin{enumerate}
    \item  \textbf{$\Tor_1^A(M,N)=0$ for all finitely-generated modules $N$.}
    We use induction on the number of generators $n$ of $N$.
    First, if $n=1$, then $N\simeq A/I$ for some ideal $I$,
    so the result is clear.
    Now, assume that the result is true for $n-1$ generators.
    Let $N$ be a module with $n$ generators
    $e_1,\dots,e_n$. Then, we have the exact sequence
    $$0\longto A[e_n] \longto N \longto N/A[e_n]\longto 0,$$
    which gives long exact sequence
    (by \Cref{thm:tor_properties} (b))
    $$\dots\longto \Tor_1^A(M,A[e_n]) \longto \Tor_1^A(M,N) \longto \Tor_1^A(M,N/A[e_n])
    \longto \dots.$$
    Since $N/A[x]$ is generated by $e_1,\dots,e_{n-1}$,
    we have by induction hypothesis that
    $\Tor_1^A(M, A[e_n]) = \Tor_1^A(M, N/A[e_n]) = 0$,
    so $\Tor_1^A(M,N) = 0$.
    \item \textbf{$\Tor_1^A(M,N)=0$ for all modules $N$.}
    Note that $N$ is a colimit of finitely-generated modules
    $\lim_{i\in I} N_i$ (by \Cref{prop:direct_limit_fingen}).
    We know that from Step 1. that $\Tor_1^A(M,N_i)=0$, so 
    by \Cref{thm:tor_properties} (d), we get that
    $$\Tor_1^A(M,N) = 
    \Tor_1^A \left(M, \varinjlim_{i\in I} N_i \right)
    = \varinjlim_{i\in I} \Tor_1^A(M,N_i)
    = \varinjlim_{i\in I} 0 = 0.$$
    \item \textbf{$M$ is flat.}
    Let $\phi : L\to N$ be injective.
    Then, the exact sequence $0\to L\to N\to N/L\to 0$
    gives  the long exact sequence
    $$0 = \Tor_1^A(N/L, M)
    \longto L\otimes M \longto N\otimes M
    \longto (N/L)\otimes M\longto 0,$$
    giving that $L\otimes M\to N\otimes M$ is injective.
    \qedhere
    \end{enumerate}
\end{proof}
How useful are these conditions? It turns out the last two have a natural interpretation.
\begin{theorem}
    $\Tor_1^A(M,A/I) = 0$ if and only if $\varphi:I\otimes_A M\to M$ given by $i\otimes m\mapsto im$ is injective.
\end{theorem}
\begin{proof}
    Consider the short exact sequence $0\to I\to A\to A/I\to 0$. Since $A\otimes M=M$ and $\Tor_1^A(A,M)=0$ (free modules are flat), we get a long exact sequence \[0\longto \Tor_1^A(A/I,M)\stackrel\alpha\longto I\otimes M\stackrel\beta\longto M\longto (A/I)\otimes M\longto 0.\] Since $\alpha$ is injective, its image is $0$ if and only if $\Tor_1^A(A/I,M)=0$. On the other hand, $\Im \alpha = \Ker \beta$, and $\Ker \beta=0$ if and only if $\beta: I\otimes M\to M$ is injective. Thus, the two are equivalent.
\end{proof}
We can now do a nontrivial example of checking 
that a morphism is flat.
\begin{example}
    Consider a ring morphism
    \begin{align*}
    \phi : \CC[z] &\to \frac{\CC[x,y]}{(xy)} \\
    z &\mapsto x+y
    \end{align*}
    (Geometrically, this corresponds to mapping point
    $(x,y)$ from variety $xy=0$ to point $z=x+y$.)
    \begin{center}
    \begin{asy}
    size(5cm,0);
    draw((-1,0)--(1,0), red+linewidth(1));
    draw((0,-1)--(0,1), red+linewidth(1));
    draw((-1,-1)--(1,1), blue+linewidth(0.8));
    for(int i=-4; i<=4; ++i){
        dot((i/5,0), red);
        dot((0,i/5),red);
        draw((i/5,0)--(i/10,i/10), gray, EndArrow(size=0.15cm));
        draw((0,i/5)--(i/10,i/10), gray, EndArrow(size=0.15cm));
    }
    \end{asy}
    \end{center}
    We claim that this makes $\CC[x,y]/(xy)$ 
    a flat $\CC[z]$-algebra.
    To prove this, note that since flatness is a local
    property, it suffices to show that
    $$M = (\CC[x,y]/(xy))_{(z-a)} \text{ is a flat }
    \CC[z]_{(x-a)}\text{-algebra is flat for all } a\in\CC.$$
    (Be careful, $M$ only allows denominator of the form
    $x+y-b$ for all $b\neq a$.)
    
    To check this, we use ideal-theoretic criterion.
    The good news is $A=\CC[z]_{(z-a)}$ has not many ideals:
    $(z-a)^n$!
    Then, it suffices to show that $\Tor_1^A(M, A/(z-a)^n) = 0$.
    To do this, we consider the free resolution
    \begin{center}
    \begin{tikzcd}[column sep = 1.5cm]
        A \arrow[r,"\cdot(z-a)^n"] & A 
    \end{tikzcd}
    \end{center}
    Tensor with $M$:
    \begin{center}
    \begin{tikzcd}[column sep = 1.5cm]
        M \arrow[r,"\cdot(x+y-a)^n"] & M
    \end{tikzcd}
    \end{center}
    Since $x+y-b$ is not a zero divisor in
    $\CC[x,y]/(xy)$ for all $b\in\CC$,
    it follows that this map is injective, so we are done.
\end{example}
\subsection{Flatness over Local Rings}
In the case of local rings, we can say a lot more about flatness.
\begin{theorem}[Flat Modules over Local Rings are Free]
    \label{thm:flat_local_ring}
    Let $A$ be a local ring with maximal ideal $\frakm$, and let $M$ be a finitely-presented $A$-module. The following are equivalent:
    \begin{alphenum}
        \item $M$ is free.
        \item $M$ is flat.
        \item $\Tor_1^A(A/\frakm, M)=0$.
    \end{alphenum}
\end{theorem}
\begin{proof}
    Note $\text{(a)}\Rightarrow\text{(b)}$ is trivial, and $\text{(b)}\Rightarrow\text{(c)}$ was proved in \Cref{thm:ideal_flat}.
    Thus, we prove that (c) implies (a).

    Suppose $\Tor_1^A(A/\frakm, M)=0$. Observe that $M/\mathfrak mM$ is a module over the field $A/\mathfrak m$, so it is free. As it is also finitely generated, we can find $m_1$, $m_2$, $\dots$, $m_n\in M$ such that their images $\overline{m_1}$, $\dots$, $\overline{m_n}$ in $M/\mathfrak mM$ form a basis. By Nakayama (\Cref{cor:nakayama_generation}), 
    $m_1$, $\dots$, $m_n$ generate $M$.

    Now define a free $A$-module $F$ with basis $\{f_1, \dots, f_n\}$, and consider the surjective map $\beta: F\to M$ which maps $f_i\mapsto m_i$ for all $i$. Let $K=\Ker \beta$, with the embedding $\alpha:K\to F$. We have a short exact sequence \[0\longto K\stackrel\alpha\longto F\stackrel\beta\longto M\longto 0.\] Note that for any module $N$, $N\otimes A/\frakm = N/\frakm N$. Thus, by tensoring the sequence with $A/\frakm$, we have an exact sequence 
    \[0 = \Tor_1^A(M, A/\frakm) \longto K/\frakm K\stackrel{\alpha_*}\longto F/\frakm F\stackrel{\beta_*}\longto M/\frakm M\longto 0.\] 
    Let $\overline{f_i}$ denote the image of $f_i$ in $F/\frakm F$. Then, $\beta_*$ maps $\overline{f_i}\mapsto \overline{m_i}$, so it maps the basis $\{\overline{f_1},\dots, \overline{f_n}\}$ of $F/\frakm F$ to the basis $\{\overline{m_1},\dots, \overline{m_n}\}$ of $M/\frakm M$. Hence, $\beta_*$ is an isomorphism. By the exactness of the sequence above, $\alpha_*$ is injective and has image $0$, so $K/\frakm K=0$. Since $M$ is finitely presented, $K$ is finitely generated
    \footnote{This is a bit subtle because by definition,
    $M$ is finitely-presented means that there exists 
    a finite presentation $K\to F\to M\to 0$
    where $F$ is finite and free and $K$ is finitely-generated.
    However, in this theorem, we are choosing specific $F$
    and $K$. Fortunately, it's true (but not obvious)
    that the definition does not depend on the choice of $F$.
    We leave this fact as an exercise to the reader.
    
    Alternatively, if you don't want to think about this issue,
    assume that $A$ is Noetherian.}, 
    so by Nakayama, $K=0$. It follows that $M\simeq F$ is free.
\end{proof}
\begin{corollary}[Equivalence of Flat and Locally Free]
    \label{cor:flat_locally_free}
    Let $A$ be a ring and $M$ be a finitely-presented $A$-module. Then, $M$ is flat if and only if $M$ is locally free (i.e., for every prime ideal $\frakp$, $M_\frakp$ is free).
\end{corollary}
\begin{proof}
    If $M$ is flat, then for all prime $\frakp$, $M_\frakp$ is a flat $A_\frakp$-module and hence free.

    If $M_\frakp$ is free for all prime $\frakp$, then $M_\frakm$ is free and hence flat for all maximal $\frakm$. By \Cref{prop:local_flat}, $M$ is flat.
\end{proof}
\subsection{(Bonus) Local Criteria of Flatness}
The above result is nice, but it assumes that $M$
is finitely-presented. 
Most morphisms of variety is not finite.
Therefore, we give a more powerful condition 
that allows us to work on non-finitely-generated ring as well.
\begin{theorem}[Local criteria of flatness]
Let $A$ and $B$ be Noetherian local rings.
Let $\frakm$ and $\frakn$ be maximal ideals of 
$A$ and $B$, respectively.
Let $\phi : A\to B$ be a \vocab{local morphism},
i.e., $\phi^{-1}\frakn = \frakm$.

Let $M$ be a finitely-generated $B$-module.
If $\Tor_1^A(M,A/\frakm) = 0$,
then $M$ is flat over $A$.
\end{theorem}
\begin{proof}
    It suffices to show that $\Tor_1^A(M,A/I)=0$
    for all ideals $I\subset A$. 
    If this is not true, take $I$ to be the maximal
    among all counterexamples
    (we can do this since $A$ is Noetherian).
    Pick $f\in\frakm\setminus I$. We have two cases. 
    \begin{itemize}
	\item \textbf{If $f$ is not a zero divisor in $A/I$},
    then we get a short exact sequence
    $$0\longto A/I\stackrel{f}\longto A/I
    \longto A/((f)+I)\longto 0.$$
    The ideal $(f)+I$ strictly contains $I$,
    so by maximality, we have $\Tor_1^A(M,A/((f)+I)) = 0$.
    Therefore, we get a surjective map
    $$\Tor_1^A(M,A/I) \stackrel{\cdot f}\longto \Tor_1^A(M,A/I).$$
    To avoid clutter, let $N=\Tor_1^A(M,A/I)$,
    which has the structure of $B$-module 
    because $M\otimes \bullet$ has the structure
    of $B$-module,
    and from above, we get a surjective map
    $N\to N$ by multiplication by $\phi(f)$.
    Then, $N=(\phi(f))N$, so by Nakayama's lemma,
    there exists $x\in (\phi(f))$ such that 
    $n = x\phi(f)n$ for all $n\in N$.
    However, from $\phi^{-1}\frakn = \frakm$, we get that
    $\phi(f)\in \frakn$. Thus, $(1-x\phi(f))$ is a unit,
    so $n=0$ for all $n\in N$, so $N=0$.
    Thus, $\Tor_1^A(M,A/I)=0$.
    \item \textbf{If $f$ is a zero divisor in $A/I$,}
    then let $J = \Ann_{A/I}(f)$ (recalling that this is the 
    set of elements $a$ such that $af=0$ 
    \Cref{def:annihilator}), so $J$ strictly contains $I$.
    Then, we have an exact sequence
    $$0\longto A/J\longto A/I\longto A/((f)+I) \longto 0.$$
    Again, by maximality assumption,
    $\Tor_1^A(M,A/J)=0$ and $\Tor_1^A(M,A/((f)+I))=0$.
    which implies by long exact sequence that 
    $\Tor_1^A(M,A/I)=0$.
    This is a contradiction. \qedhere
    \end{itemize}
\end{proof}
\section{Going-Down Theorem, Flat Version}
\label{sec:going_down}
In \Cref{thm:going_down}, we proved a version of going down for integral extensions of integrally-closed domains. In this section, we will prove another version of the statement which instead assumes flatness.
\begin{theorem}[Going-Down Theorem, Flat Version]
    \label{thm:going_down_flat}
    Given a morphism $\phi : A\to B$, 
    we say that prime ideal $\qq\subset B$
    \vocab{lies over} $\pp\subset A$
    if $\phi^{-1}\qq = \pp$.
    
    Suppose that $\phi  :A\to B$ is a flat morphism.
    Given that $\mathfrak p\subset\mathfrak p'\subset A$
    be prime ideals and $\mathfrak q'$ lies over $\mathfrak p'$,
    then there exists a prime ideal $\mathfrak q\subset\mathfrak q'$
    lying over $\mathfrak p$.
    \begin{center}
    \begin{tikzcd}
        {\color{gray}\qq}  \arrow[d, no head, dashed] \arrow[r,"\subset",phantom]
        & \qq'  \arrow[d, no head] \arrow[r, "\subset",phantom]
        & B \arrow[d, no head] \\
        \pp   \arrow[r,"\subset",phantom]
        & \pp'  \arrow[r, "\subset",phantom]
        & A 
    \end{tikzcd}
    \end{center}
\end{theorem}
Before we prove this, we will need a new definition.
\subsection{Faithful Flatness}
One way of stating the definition of flatness is as follows: a module $M$ is flat if $L\to N\to P$ being exact implies $L\otimes M\to N\otimes M\to P\otimes M$ is exact.
\begin{definition}
    A module $M$ is \vocab{faithfully flat} if the following is true: $L\to N\to P$ is exact \textbf{if and only if} $L\otimes M\to N\otimes M\to P\otimes M$ is exact.

    Similarly, a morphism of rings $\varphi : A\to B$
    is \vocab{faithfully flat} if it makes $B$
    a faithfully flat $A$-algebra.
\end{definition}
\begin{example}
    Finitely-generated free modules are faithfully flat, as $N\otimes_A A^{\oplus n} \cong N^{\oplus n}$. (In fact, all free modules are faithfully flat.)

    An example of a flat module which is not faithfully flat is $\QQ$ as a $\ZZ$-module. Indeed, by \Cref{cor:localization_flat}, $\QQ = (\ZZ\setminus \{0\})^{-1}\ZZ$ is flat. Further note $(\ZZ/n\ZZ)\otimes \QQ = 0$, as $m\otimes q = (nm)\otimes (q/n) = 0$ for all $m$ and $q$. Thus, while \[\ZZ/n\ZZ\stackrel{\cdot 1}\longto \ZZ/ n\ZZ\stackrel{\cdot 1}\longto \ZZ/n\ZZ\] 
    is not exact, its tensor with $\QQ$ is $0\to 0\to 0$, which is exact.
\end{example}
\begin{definition}[Residue Field]
    Let $A$ be a ring and $\pp\subseteq A$ be a prime ideal. The \vocab{residue field} is defined as
    \notation{$\kappa(\pp)$}{residue field}
    \[\kappa(\pp) := \Frac(A/\pp).\]
\end{definition}
\begin{theorem}
    \label{thm:faithful_flat_conditions}
    Let $M$ be a flat $A$-module. The following are equivalent;
    \begin{alphenum}
        \item $M$ is faithfully flat.
        \item If $N\otimes_A M=0$, then $N=0$.
        \item For all prime ideals $\pp\subseteq A$, $\kappa(\pp)\otimes_A M \neq 0$.
        \item For all maximal ideals $\frakm\subseteq A$, $\kappa(\frakm)\otimes_A M \neq 0$.
    \end{alphenum}
\end{theorem}
\begin{proof}
    It is clear that $\text{(b)}\Rightarrow\text{(c)}\Rightarrow\text{(d)}$.

    \textbf{(a) $\Rightarrow$ (b).} Suppose $M$ is faithfully flat. If $M\otimes N=0$, then $0\to M\otimes_A 0\to M\otimes_A N\to M\otimes_A 0\to 0$ is exact, so $0\to 0\to N\to 0\to 0$ is exact. Thus, $N=0$.

    \textbf{(b) $\Rightarrow$ (a).} Suppose $N\otimes_A M=0$ only when $N=0$, and suppose for modules $L$, $N$, and $P$, the tensored sequence \[L\otimes_A M\stackrel{\alpha_*}\longto N\otimes_A M\stackrel{\beta_*}\longto P\otimes_A M\] is exact. Then, $\Im(\beta\alpha)\otimes_A M = \Im(\beta_*\alpha_*)=0$, so $\Im(\beta\alpha) = 0$, i.e., $\Im \alpha\subseteq \Ker \beta$. By flatness, \[\frac{\Ker\beta}{\Im\alpha} \otimes_A M = \frac{\Ker\beta\otimes_A M}{\Ker\alpha\otimes_A M} = \frac{\Ker\beta_*}{\Im\alpha_*} = 0,\] so $\Ker\beta / \Im \alpha = 0$, i.e., $L\to N\to P$ is exact. Thus, $M$ is faithfully flat.

    \textbf{(d) $\Rightarrow$ (b).} Suppose $M\otimes_A \kappa(\frakm) \neq 0$ for all maximal $\frakm$. Consider any nonzero module $N$. Let $x\in N$ be a nonzero element. Observe $A/\Ann(x)$ embeds into $N$ under the map, so $(A/\Ann(x)) \otimes_A M$ embeds into $N\otimes_A M$.
    
    Let $\frakm$ be any maximal ideal containing $\Ann(x)$. Then, $(A/\Ann(x))\otimes_A M$ surjects onto $(A/\frakm)\otimes_A M$. The latter is nonzero by assumption, so the former is also nonzero. It embeds into $N\otimes_A M$, so $N\otimes_A M\neq 0$, as desired.
\end{proof}
The module $M\otimes_A \kappa(\pp)$ has a geometric
significance: it is a \vocab{fiber}.
To see this, we first note that since localization
commutes with quotient, we get 
$\kappa(\pp) \simeq A_\pp/\pp A_\pp$, so we have the 
exact sequence
$$\pp A_\pp \longto A_\pp \longto \kappa(\pp) \longto 0,$$
so tensoring gives
$$\pp M_\pp \longto M_\pp \longto \kappa(\pp)\otimes_A M\longto 0,$$
which implies that
\begin{equation}
\label{eq:fiber}
\kappa(\pp)\otimes_A M \simeq M_\pp / \pp M_\pp.
\end{equation}

Now, suppose that $M=B$ is an $A$-algebra.
Then, $\kappa(\pp) \otimes_A B \simeq B_\pp/\pp B_\pp$
as $A$-algebra as well.
We claim the following:
\begin{proposition}[Points in fiber indeed corresponds to preimage]
    \label{prop:fiber}
    If $\phi : A\to B$ be a morphism
    (so $B$ is an $A$-algebra), then 
    there is a bijective correspondence between
    $$\{\text{prime ideals of }\kappa(\pp)\otimes_A B\}
    \longleftrightarrow 
    \{\text{prime ideals } \qq\subset B 
    \text{ such that } \phi^{-1}\qq = \pp\}.$$
\end{proposition}
\begin{proof}
    Recall \eqref{eq:fiber} that $\kappa(\pp)\otimes_A B 
    = B_\pp/\pp B_\pp$.
    First, note that $B_\pp$ is a localization of $B$
    with multiplicative set $\phi(A/\pp)$.
    Thus, every prime ideal of $B_\pp$ (and hence $B_\pp/\pp B_\pp$)
    is extended from a prime ideal $\qq\in B$.
    Observe the following:
    \begin{itemize}
    \item In order to get a prime ideal in $B_\pp$,
    we must have $\qq\cap \phi(A\setminus\pp) = \emptyset$,
    This is equivalent to $\phi^{-1}\qq \subseteq \pp$
    (which we left as an easy exercise to check).
    \item In order to get a prime ideal in $B_\pp/\pp B_\pp$,
    we must have $\qq A_\pp \supseteq \phi(\pp)A_\pp$,
    which means $\qq\supseteq\phi(\pp)$
    or equivalently $\phi^{-1}\qq \supseteq \pp$.
    \end{itemize}
    Combining both items gives $\qq$ extends to 
    a prime ideal in $B_\pp/\pp B_\pp$ if and only if
    $\phi^{-1}\qq=\pp$.
\end{proof}
In particular, $\Spec(\kappa(\pp)\otimes_A B)
\simeq \phi^{-1}(\{\pp\})$.
This is why $\kappa(\pp)\otimes_A B$ is 
an algebraic realization of the fiber of
$\phi^\# : \Spec B\to\Spec A$ at $\pp$.

\subsection{Proof of Going Down}
\begin{corollary}[Faithfully flat implies surjective]
    \label{cor:faithfully_flat_contraction}
    If $\varphi:A\to B$ is faithfully flat (i.e., makes $B$ a faithfully flat $A$-module), then $\varphi^\#: \Spec B\to \Spec A$ is surjective.
\end{corollary}
\begin{proof}
    Let $\pp$ be any prime ideal of $A$.
    From \Cref{thm:faithful_flat_conditions},
    $\kappa(\pp) \otimes_A B\neq 0$, so there exists
    at least one prime ideal, which by \Cref{prop:fiber}
    must then correspond
    to prime ideal in $\Spec B$ that contracts to $\pp$.
\end{proof}
\begin{corollary}
    \label{cor:faithfully_flat_lying}
    If $\varphi:A\to B$ is flat and $\qq\in \Spec B$ lies over $\pp\in \Spec A$, then $B_\qq$ is faithfully flat over $A_\pp$.
\end{corollary}
\begin{proof}
    First, we show flatness. 
    Let $B_\pp = (A\setminus\pp)^{-1} B$ be an $A_\pp$-module
    (warning: it is not an algebra).
    Then, by \Cref{cor:localization_flat}, we get that
    $B_\pp$ is a flat $A_\pp$-module.
    Localizing $B_\pp$ further gives that $B_\qq$
    is a flat $B_\pp$-module, hence a flat $A_\pp$-module
    by forgetting scalar.

    Now, we show faithful flatness.
    Since $A_\pp$ is local with maximal ideal $\pp A_\pp$, which $\qq B_\qq\subseteq B_\qq$ lies over, $\pp A_\pp$ is in the image of $\varphi^\#: \Spec B_\qq\to\Spec A_\pp$. Thus, $B_\qq$ is faithfully flat over $A_\pp$.
\end{proof}
\begin{proof}[Proof of Going Down \Cref{thm:going_down_flat}]
    Suppose $B$ is a flat $A$-module, and define $\pp$, $\pp'$, $\qq'$ as in the statement. By \Cref{cor:faithfully_flat_lying}, $B_{\qq'}$ is faithfully flat over $A_{\pp'}$. By \Cref{cor:faithfully_flat_contraction}, some prime in $B_{\qq'}$ lies over $\pp A_{\pp'}$. Lifting this prime back to $B$, we get a prime that lies over $\pp$.
\end{proof}

\section{Flatness Implies Constant-Dimension Fibers}
\label{sec:flat_const_dim}
We now prove an important result that flat morphisms
have constant dimension of fibers.
\begin{theorem}[Flatness implies constant dimension of fibers]
    \label{thm:flat_dim_fiber}
    If $A$ and $B$ are Noetherian local rings with maximal ideals
    $\frakm$ and $\frakn$ and $\varphi : A\to B$
    is a ring morphism such that $\varphi^{-1}\frakn = \frakm$.
    Then,
    $$\dim B \leq \dim A + \dim B/\frakm B.$$
    Moreover, if $\varphi$ is flat, then the equality holds.
\end{theorem}
Before proving this theorem, let us consider geometric interpretation.
Let $\varphi : A\to B$ be a flat morphism,
$\qq$ be a prime ideal of $B$, and $\pp = \varphi^{-1}A$.
Then, $\varphi$ induces a local morphism
$\varphi_* : A_\pp \to B_\qq$. Then, we have
\begin{equation}
\label{eq:dim_fiber}
\dim B_\qq \leq \dim A_\pp + \dim B_\qq / \pp B_\qq.
\end{equation}
Note that by localization in stages,
$$B_\qq / \pp B_\qq = (B_\pp / \pp B_\pp)_\qq 
\stackrel{\eqref{eq:fiber}}= (B \otimes_A \kappa(\pp))_\qq$$
(where $B_\pp = (\phi(A\setminus\pp))^{-1}B$).
The $A$-algebra $B\otimes_A \kappa(\pp)$ is the \vocab{fiber} 
of the morphism $\varphi^{\#} : \Spec B\to \Spec A$ 
at $\pp$ (i.e., $\Spec(B\otimes_A \kappa(\pp)) = 
(\varphi^\#)^{-1}(\{\pp\})$, which follows from 
\Cref{prop:fiber}). 
This theorem says that for flat morphisms,
the local dimension of fiber at $\qq$
is the difference of the local dimension of $B$ at $\qq$
and of $A$ at $\pp$.
\begin{example}
    Consider the morphism
    $$\phi : k[x] \to \frac{k[x,y]}{(xy)}$$
    mentioned in the introduction. Let $A=k[x]$,
    $B=k[x,y]/(xy)$, $\qq = (x,y)\subseteq B$,
    so $\pp  = \phi^{-1}\qq = (x) \subseteq A$.
    Then, using the fact that localization commutes
    with quotient, we compute
    $$\frac{B_\qq}{\pp B_\qq} 
    = \frac{B_{(x,y)}}{(x)B_{(x,y)}}
    = \left(\frac{B}{(x)}\right)_{(x,y)}
    = k[y]_{(y)},$$
    which has dimension $1$. 
    This correspond to the geometric intuition
    that inverse image of $\{0\}$ is the $y$-axis.
    
    On the other hand, $\dim A_\pp = \dim B_\qq = 1$
    by \Cref{prop:local_ring_slice}.
    This violates \eqref{eq:dim_fiber}.
\end{example}
\begin{proof}[Proof of \Cref{thm:flat_dim_fiber}]
    Let $d=\dim A$ and $e=\dim B/\mm B$. 
    From primary ideal dimension, there exists an $\mm$-primary ideal $(x_1,\dots, x_n)\subseteq A$, i.e., there exists $n_1$ for which $\mm^{n_1} \subseteq (x_1,\dots, x_n)\subseteq \mm$ (this utilizes $A$ is Noetherian).
    Similarly, there exists $n_2$ and an $\frakn/\mm B$-primary ideal $(y_1,\dots, y_e)\subseteq \frakn$ containing $\mm B$ such that 
    \[\left(\frac{\mathfrak n}{\mathfrak mB}\right)^{n_2} \subseteq \frac{(y_1,\dots, y_e)}{\mathfrak mB}\subseteq \frac{\mathfrak n}{\mathfrak mB}.\] Since $B/\mm B$ is local with maximal ideal $\frakn /\mm B$, we know $\frakn^{n_2} \subseteq (y_1,\dots, y_e) + \mm B$. Thus, 
    \begin{align*}
    \frakn^{n_1n_2} &\subseteq ((y_1,\dots, y_e) + \mm B)^{n_1} 
    \\ 
    &\subseteq (y_1,\dots, y_e) + (x_1,\dots, x_d)B \\
    &=(y_1,\dots, y_e,\varphi(x_1),\dots, \varphi(x_d))  \\
    &\subseteq \frakn,
    \end{align*}
    so there exists an $\nn$-primary ideal with $d+e$ generators.
    We conclude $\dim B\le d+e$, the desired inequality.
    
    Now suppose $\varphi$ is flat. Given a chain of prime ideals
    \begin{align*}
        \pp_0\subsetneq \pp_1\subsetneq\dots\subsetneq \pp_d 
        = \frakm
        \ &\subseteq \ A \\
        \frac{\qq_0}{\mm B}\subsetneq \frac{\qq_1}{\mm B}\subsetneq\dots\subsetneq \frac{\qq_e}{\mm B} 
        \ &\subseteq \ \frac{B}{\mm B}
    \end{align*}
    (where $\qq_0\supseteq \mm B$), one construct 
    a chain of prime ideals in $B$ as follows.
    Note that $\mm \supseteq \varphi^{-1}\qq_0\supseteq \varphi^{-1}\mm B\supseteq \mm$, so
    $\varphi^{-1}\qq_0=\mm$, so $\qq_0$ lies above $\pp_d=\mm$,
    so by going down theorem \Cref{thm:going_down_flat}, 
    there exists prime ideals 
    $\overline{\pp}_i$ such that $\varphi^{-1}\overline\pp_i
    = \pp_i$.
    The whole construction can be summarized
    in the diagram below.
    \begin{center}
    \begin{tikzcd}[column sep=0.7cm]
    B/\mm B: & & & & 
    & \dfrac{\qq_0}{\mm B}
    \arrow[r, "\subsetneq", phantom]\arrow[d, no head]
    &\dfrac{\qq_1}{\mm B}
    \arrow[r, "\subsetneq", phantom]\arrow[d, no head]
    & \dots 
    \arrow[r, "\subsetneq", phantom]
    & \dfrac{\qq_e}{\mm B}
    \arrow[d, no head]
    \\
    B:
    & {\overline\pp_0}  
    \arrow[r, "\subsetneq", phantom] \arrow[d, no head, dashed]
    & {\overline\pp_1} 
    \arrow[r, "\subsetneq", phantom] \arrow[d, no head, dashed]
    & \dots \arrow[r, "\subsetneq", phantom]
    & {\overline\pp_{d-1}} 
    \arrow[r, "\subsetneq", phantom] \arrow[d, no head, dashed]
    & \qq_0 \arrow[r, "\subsetneq", phantom]
    \arrow[d, no head]
    & \qq_1 \arrow[r, "\subsetneq", phantom]
    & \dots \arrow[r, "\subsetneq", phantom]
    & \qq_e \\
    A: & \pp_0 \arrow[r, "\subsetneq", phantom]
    & \pp_1 \arrow[r, "\subsetneq", phantom]
    & \dots \arrow[r, "\subsetneq", phantom]
    & \pp_{d-1} \arrow[r, "\subsetneq", phantom]
    & \pp_d = \mm
    \end{tikzcd}
    \end{center}
    Now, our chain of prime ideal is 
    $\overline\pp_0\subsetneq \overline\pp_1 
    \subsetneq\dots \subsetneq \overline\pp_{d-1} 
    \subsetneq \qq_0 \subsetneq \qq_1 
    \subsetneq\dots\subsetneq \qq_e$.
\end{proof}
As an application, 
we now give the deferred proof of \Cref{prop:dim_poly}.
\begin{corollary}[Dimension of polynomial ring]
    \label{cor:dim_poly}
    Let $A$ be a Noetherian ring. Then, $\dim A[x] = \dim A+1$.
\end{corollary}
\begin{proof}
    First, we note that the morphism $A\to A[x]$ is flat
    because $A[x]$ is a free $A$-module.
    Let $\qq$ be a prime ideal of $A[x]$, and let 
    $\pp$ be its contraction to $A$. By the previous corollary,
    we get that
    $$\dim A[x]_\qq = \dim A_\pp + \dim(A[x]_\qq / \pp A[x]_\qq).$$
    However, we note that since localization commutes 
    with everything,
    $$A[x]_\qq / \pp A[x]_\qq = (A[x] / \pp A[x])_{\qq} 
    = ((A/\pp)[x])_{\qq} = \kappa(\pp)[x]_\qq,$$
    which has dimension $1$. Therefore,
    $$\dim A[x]_\qq = \dim A_\pp + 1$$
    for all prime ideal $\qq$. Taking maximum across
    all prime ideal $\qq$ (noting that $\Spec A[x] \to \Spec A$
    is surjective because $\pp[x]$ maps to $\pp$),
    we get that $\dim A[x]=\dim A+1$.
\end{proof}
\section{Problems}
\begin{problem}
\label{prob:not_flat_example}
Prove that $\CC[x,y]/(xy)$ is not a flat $\CC[x]$-algebra.
\end{problem}
\begin{problem}
Let $M$ be a module over ring $k[t]/(t^2)$.
Prove that $M$ is flat if and only if the multiplication by $t$ map
$M/tM \to tM$ is an isomorphism.
\end{problem}
\begin{problem}
Let $A=\CC[x,y]/(xy)$. Compute $\Tor_n^A(A/(x), A/(x))$.
Your answer should depend on whether $n=0$, $n$ is odd,
or $n$ is even and positive.
\end{problem}
\begin{problem}
Let $A$ be a regular local ring of dimension $n$ with maximal ideal $\frakm$
and residue field $k$.
From the proof of \Cref{prob:cohen_macaulay} (c), 
there exists a sequence $x_1,x_2,\dots,x_n\in \mm$ 
such that $\mm = (x_1,\dots,x_n)$ and 
$x_i$ is not a zero-divisor of $A/(x_1,\dots,x_{i-1})$.
\begin{alphenum}
\item Prove that for any $A$-module $M$, we have 
$\Tor^A_i(M,k) = 0$ for all $i>n$.

(Hint: There are at least two ways to do this.
One is to use $0\to A/(x_1,\dots,x_i) \to A/(x_1,\dots,x_i)
\to A/(x_1,\dots,x_{i+1})\to 0$.
The other is to explicitly find a free resolution of $k$)
\item Prove that any $A$-module $M$ 
has a free resolution $F_n\to \dots \to F_1\to F_0$
of length at most $n$.\footnote{The following converse holds:
if $A$ is Noetherian local ring such that the residue field $k$
has a finite free resolution, then $A$ is regular.
See \cite[Theorem 19.12]{eisenbud} for the proof.}

(Hint: start with any free resolution $\dots \to F_1\to F_0$.
Truncate it to a resolution $K\to F_{n-1}\to\dots\to F_1\to F_0$
where $K=\Im(F_n\to F_{n-1})$
and use \Cref{thm:flat_local_ring} to show that $K$ is free.)
\item For each $n$, given an example showing that the bound $n$ in (b) cannot be improved.
\end{alphenum}
\end{problem}